\documentclass[12pt,a4paper]{article}

\usepackage[utf8]{inputenc}
\usepackage[T1]{fontenc}
\usepackage{polski}
\usepackage[polish]{babel}
\usepackage{graphicx}
\usepackage{hyperref}
\usepackage[margin=2.5cm]{geometry}
\usepackage{listings}
\usepackage{xcolor}
\usepackage{amsmath}
\usepackage{enumitem}
\usepackage{tikz}
\usepackage{float}
\usepackage{titlesec}
\usepackage{colortbl}
\usepackage{array}
\usepackage{pifont}
\usepackage{booktabs}
\usepackage{longtable}

\usetikzlibrary{arrows.meta,positioning,shapes.geometric}

\newcommand{\cmark}{\ding{51}}
\newcommand{\xmark}{\ding{55}}

\definecolor{lightgray}{gray}{0.9}
\definecolor{navyblue}{RGB}{0,31,91}
\definecolor{headertext}{RGB}{255,255,255}
\definecolor{codegreen}{rgb}{0,0.6,0}
\definecolor{codegray}{rgb}{0.5,0.5,0.5}
\definecolor{codepurple}{rgb}{0.58,0,0.82}
\definecolor{backcolour}{rgb}{0.95,0.95,0.95}

\newcolumntype{H}{>{\columncolor{lightgray}\bfseries}l}

\titleclass{\subsubsubsection}{straight}[\subsubsection]
\newcounter{subsubsubsection}[subsubsection]
\renewcommand\thesubsubsubsection{\thesubsubsection.\arabic{subsubsubsection}}
\titleformat{\subsubsubsection}
  {\normalfont\normalsize\bfseries}{\thesubsubsubsection}{1em}{}
\titlespacing*{\subsubsubsection}{0pt}{3.25ex plus 1ex minus .2ex}{1.5ex plus .2ex}
\setcounter{tocdepth}{4}

\lstdefinestyle{mystyle}{
    backgroundcolor=\color{backcolour},
    commentstyle=\color{codegreen},
    keywordstyle=\color{blue},
    numberstyle=\tiny\color{codegray},
    stringstyle=\color{codepurple},
    basicstyle=\ttfamily\scriptsize,
    breakatwhitespace=false,
    breaklines=true,
    captionpos=b,
    keepspaces=true,
    numbers=left,
    numbersep=5pt,
    showspaces=false,
    showstringspaces=false,
    showtabs=false,
    tabsize=2
}
\lstset{style=mystyle}

\hypersetup{
    colorlinks=true,
    linkcolor=navyblue,
    urlcolor=navyblue,
    citecolor=navyblue
}

\begin{document}

\begin{titlepage}
\begin{center}
\large
    {\noindent Collegium Witelona Uczelnia Państwowa w Legnicy}\\
    {\noindent Wydział Nauk Technicznych i Ekonomicznych}\\
    {\noindent Kierunek: Informatyka}\\[2cm]
    \vspace{1cm}
    {\Huge\textbf{Siam Skirmish}}\\[0.4cm]
    {\Large\textbf{Dokumentacja projektu gry turowej RPG}}\\[0.8cm]
    {\large Projekt wykonany w silniku Unity}\\[3cm]
    {\large\textbf{Autor}}\\
    Gabriela Grabarska, nr indeksu: 43840\\[2.5cm]
\end{center}

\begin{flushright}
    \begin{tabular}{l}
        Prowadzący przedmiot\\
        mgr inż. Zdzisław Pawelec
    \end{tabular}
\end{flushright}

\vfill
\begin{center}
    {\noindent Legnica, 2026}
\end{center}
\end{titlepage}

\tableofcontents
\newpage

% ============================================================
\section{Cel projektu i powód jego powstania}
% ============================================================

\subsection{Problem, który rozwiązujemy}

Projekt \textit{Siam Skirmish} jest turową grą taktyczną RPG, w której gracz kieruje drużyną postaci walczących na planszy kafelkowej z przeciwnikami sterowanymi przez komputer. Głównym problemem projektowym było stworzenie przejrzystego systemu walki turowej, który łączy prostą obsługę, czytelną planszę, różne klasy jednostek oraz przeciwników podejmujących decyzje w sposób bardziej elastyczny niż zwykłe losowanie ruchu.

Gra realizuje klasyczny model taktycznej potyczki: jednostki poruszają się po siatce, wykonują akcje ruchu i ataku, a wynik starcia zależy od pozycji, zasięgu, klasy postaci, punktów życia i postawy bojowej.

\subsection{Dlaczego projekt jest istotny?}

W wielu prostych grach turowych sztuczna inteligencja przeciwnika działa na podstawie sztywnych instrukcji: zawsze podchodzi do najbliższego celu albo zawsze atakuje, jeśli tylko może. Taki model jest łatwy do implementacji, ale prowadzi do przewidywalnej rozgrywki. W projekcie zastosowano logikę rozmytą, dzięki której przeciwnik może ocenić sytuację na planszy w sposób stopniowy: czy ma mało punktów życia, czy cel jest blisko, czy zagrożenie ze strony graczy jest wysokie oraz czy bardziej opłaca się przyjąć postawę defensywną, zbalansowaną czy agresywną.

\subsection{Cel projektu}

Celem projektu jest stworzenie działającej gry 2D w Unity, która:

\begin{itemize}
    \item posiada planszę kafelkową z jednostkami gracza i przeciwnika,
    \item obsługuje system tur i kolejkę postaci,
    \item pozwala graczowi wykonywać ruch i atak w wybranej kolejności,
    \item posiada przeciwników sterowanych przez AI oparte na logice rozmytej,
    \item prezentuje stan gry przez interfejs użytkownika, log walki, kolejkę tur i paski życia,
    \item rozróżnia klasy jednostek z innymi statystykami,
    \item kończy rozgrywkę sceną zwycięstwa albo porażki.
\end{itemize}

\newpage

% ============================================================
\section{Opis gry}
% ============================================================

\subsection{Gatunek i koncepcja}

\textit{Siam Skirmish} to taktyczna gra RPG 2D z widokiem z góry. Rozgrywka odbywa się na planszy zbudowanej z kafelków. Każda jednostka posiada pozycję na siatce, punkty życia, zasięg ruchu, zasięg ataku oraz obrażenia zależne od klasy.

Gracz kontroluje drużynę bohaterów, natomiast przeciwnicy są obsługiwani automatycznie przez system AI. Gra kończy się zwycięstwem, gdy wszyscy przeciwnicy zostaną pokonani, albo porażką, gdy wszystkie jednostki gracza zostaną wyeliminowane.

\subsection{Podstawowa pętla rozgrywki}

\begin{enumerate}
    \item Gra uruchamia scenę główną i generuje planszę.
    \item System losuje oraz rozmieszcza jednostki gracza i przeciwnika.
    \item Menedżer tur buduje kolejkę żywych jednostek.
    \item Aktywna jednostka wykonuje swoją turę.
    \item Gracz wybiera ruch, atak albo zakończenie tury.
    \item Przeciwnik analizuje sytuację przez AI rozmyte i wykonuje decyzję.
    \item Po każdej akcji sprawdzany jest warunek zwycięstwa albo porażki.
    \item Po zakończeniu tury aktywowana jest kolejna jednostka z kolejki.
\end{enumerate}

\subsection{Najważniejsze założenia}

\begin{longtable}{p{4cm} p{10cm}}
\toprule
\textbf{Obszar} & \textbf{Założenie projektowe} \\
\midrule
\endhead
Plansza & Gra działa na siatce kafelkowej o rozmiarze 12 na 12 pól. \\
Ruch & Jednostka może poruszyć się tylko na osiągalne, wolne pole. \\
Walka & Atak jest możliwy tylko na przeciwnika w zasięgu danej klasy. \\
Tury & W jednej turze jednostka może wykonać ruch i atak albo tylko jedną z tych akcji. \\
AI & Przeciwnik wybiera cel, postawę, miejsce ruchu i kolejność akcji. \\
UI & Interfejs pokazuje aktywną jednostkę, fazę tury, kolejkę, log walki i paski życia. \\
Koniec gry & Gra przechodzi do osobnej sceny zwycięstwa albo porażki. \\
\bottomrule
\end{longtable}

\newpage

% ============================================================
\section{Podział funkcjonalności}
% ============================================================

\subsection{Funkcje dostępne dla gracza}

\begin{longtable}{>{\bfseries}p{1cm} p{4cm} p{8cm} p{2cm}}
\toprule
ID & Funkcja & Opis & Priorytet \\
\midrule
\endhead
1 & Rozpoczęcie gry & Gracz uruchamia rozgrywkę z menu głównego. & Niezbędne \\[4pt]
2 & Wybór ruchu & Gracz wybiera pole, na które aktywna jednostka ma się przemieścić. & Niezbędne \\[4pt]
3 & Wybór ataku & Gracz wybiera przeciwnika znajdującego się w zasięgu ataku. & Niezbędne \\[4pt]
4 & Pominięcie kroku & Gracz może pominąć ruch lub atak prawym przyciskiem myszy albo przyciskiem interfejsu. & Przydatne \\[4pt]
5 & Zakończenie tury & Gracz może zakończyć turę aktywnej jednostki. & Niezbędne \\[4pt]
6 & Podgląd zasięgu & Gra podświetla pola ruchu i pola ataku. & Niezbędne \\[4pt]
7 & Podgląd kolejki tur & Interfejs pokazuje, które jednostki wykonają ruch jako następne. & Przydatne \\[4pt]
8 & Obserwacja logu walki & Gracz widzi komunikaty o zadanych obrażeniach i pokonanych jednostkach. & Przydatne \\[4pt]
9 & Podgląd punktów życia & Nad każdą jednostką widoczny jest pasek życia. & Niezbędne \\
\bottomrule
\end{longtable}

\subsection{Funkcje realizowane automatycznie przez system}

\begin{longtable}{>{\bfseries}p{1cm} p{4.5cm} p{9cm}}
\toprule
ID & Funkcja systemowa & Opis \\
\midrule
\endhead
1 & Generowanie planszy & System tworzy wizualną siatkę kafelków i obsługuje mapowanie pozycji świata na pola planszy. \\[4pt]
2 & Rozmieszczenie jednostek & System losuje wolne pola startowe po stronie gracza i przeciwnika. \\[4pt]
3 & Budowanie kolejki tur & Menedżer tur zbiera żywe jednostki i ustala kolejność wykonywania akcji. \\[4pt]
4 & Walidacja akcji & System sprawdza, czy ruch i atak są zgodne z regułami gry. \\[4pt]
5 & Obliczanie obrażeń & System dobiera najlepszy typ ataku według dystansu i modyfikuje obrażenia przez postawę. \\[4pt]
6 & Decyzje AI & Przeciwnik wybiera cel, postawę, ruch oraz sekwencję działań na podstawie logiki rozmytej. \\[4pt]
7 & Aktualizacja UI & Interfejs odświeża aktywną jednostkę, fazę, kolejkę, panel celu i log. \\[4pt]
8 & Aktualizacja pasków życia & Paski życia skalują się proporcjonalnie do aktualnych punktów życia jednostki. \\[4pt]
9 & Sprawdzanie końca gry & System wykrywa zwycięstwo albo porażkę i przełącza scenę. \\
\bottomrule
\end{longtable}

\newpage

% ============================================================
\section{Dane wejściowe i wyjściowe}
% ============================================================

\subsection{Dane wejściowe}

\begin{longtable}{p{4cm} p{5cm} p{5cm}}
\toprule
\textbf{Dane} & \textbf{Źródło} & \textbf{Zastosowanie} \\
\midrule
\endhead
Kliknięcie lewym przyciskiem myszy & \texttt{PlayerInputController} & Wybór pola ruchu albo celu ataku. \\[4pt]
Kliknięcie prawym przyciskiem myszy & \texttt{PlayerInputController} & Pominięcie aktualnego kroku tury. \\[4pt]
Parametry klas jednostek & Assety \texttt{UnitClassData} & Maksymalne HP, ruch, zasięgi i obrażenia. \\[4pt]
Pozycje startowe & \texttt{BattleSpawner} & Losowe ustawienie drużyn na wolnych polach. \\[4pt]
Stan planszy & \texttt{GridManager} & Sprawdzanie pól wolnych, blokad i zasięgów. \\[4pt]
Stan jednostek & \texttt{Unit} & Punkty życia, drużyna, postawa, pozycja i aktywność. \\[4pt]
Stan walki & \texttt{TurnManager} & Aktualna faza tury, aktywna jednostka i kolejka. \\
\bottomrule
\end{longtable}

\subsection{Dane wyjściowe}

\begin{longtable}{p{4cm} p{5cm} p{5cm}}
\toprule
\textbf{Dane} & \textbf{Forma prezentacji} & \textbf{Opis} \\
\midrule
\endhead
Plansza & Tilemapa i sprite tła & Obszar walki widoczny dla kamery. \\[4pt]
Jednostki & SpriteRenderer & Postacie gracza i przeciwników na kafelkach. \\[4pt]
Zasięg ruchu & Podświetlenie kafelków & Pola, na które jednostka może wejść. \\[4pt]
Zasięg ataku & Podświetlenie kafelków & Pola możliwego ataku. \\[4pt]
Paski życia & SpriteRenderery nad postaciami & Aktualny procent HP każdej jednostki. \\[4pt]
Log walki & Panel UI & Komunikaty o atakach, obrażeniach i eliminacjach. \\[4pt]
Kolejka tur & Panel UI & Następne jednostki w kolejności aktywacji. \\[4pt]
Wynik gry & Scena końcowa & Scena zwycięstwa lub porażki. \\
\bottomrule
\end{longtable}

\newpage

% ============================================================
\section{Architektura systemu}
% ============================================================

\subsection{Technologie}

\begin{longtable}{p{4cm} p{4cm} p{6cm}}
\toprule
\textbf{Element} & \textbf{Technologia} & \textbf{Rola w projekcie} \\
\midrule
\endhead
Silnik gry & Unity & Uruchamianie scen, obsługa obiektów, renderowanie 2D. \\[4pt]
Język programowania & C\# & Implementacja logiki gry i komponentów. \\[4pt]
Renderowanie 2D & Universal Render Pipeline 2D & Obsługa kamery, sprite'ów, tilemap i efektów 2D. \\[4pt]
Interfejs & Unity UI / UGUI & Przyciski, panele, teksty, log walki i menu. \\[4pt]
Plansza & Unity Tilemap & Reprezentacja pól i warstw podświetleń. \\[4pt]
Dane klas & ScriptableObject & Konfiguracja statystyk jednostek bez zmiany kodu. \\
\bottomrule
\end{longtable}

\subsection{Główne moduły projektu}

\begin{longtable}{p{4cm} p{5cm} p{6cm}}
\toprule
\textbf{Moduł} & \textbf{Najważniejsze skrypty} & \textbf{Odpowiedzialność} \\
\midrule
\endhead
AI & \texttt{FuzzyAiController}, \texttt{FuzzyMembership} & Decyzje przeciwników, logika rozmyta, wybór celu i sekwencji akcji. \\[4pt]
Combat & \texttt{CombatSystem} & Walidacja ataku, obliczanie obrażeń i komunikaty walki. \\[4pt]
Grid & \texttt{GridManager}, \texttt{GridCoordinateSystem} & Plansza, kafelki, współrzędne, zasięg ruchu i ataku. \\[4pt]
Input & \texttt{PlayerInputController} & Obsługa kliknięć gracza. \\[4pt]
Spawning & \texttt{BattleSpawner} & Tworzenie jednostek i rozmieszczanie ich na planszy. \\[4pt]
Turns & \texttt{TurnManager}, \texttt{ActionValidator}, \texttt{UnitQuery}, \texttt{BattleState} & Kolejka tur, stan walki, wybór akcji i przejścia między fazami. \\[4pt]
Units & \texttt{Unit}, \texttt{UnitView}, \texttt{UnitClassData}, \texttt{UnitHealthBar} & Dane jednostek, ruch wizualny, klasy postaci i paski życia. \\[4pt]
UI & \texttt{BattleUIController}, \texttt{TargetPanelUI}, \texttt{CombatLogUI}, \texttt{LogManager}, \texttt{MainMenuUI}, \texttt{VictoryUI}, \texttt{GameOverUI} & Prezentacja stanu gry i ekranów menu. \\
\bottomrule
\end{longtable}

\subsection{Diagram przepływu rozgrywki}

\begin{figure}[H]
\centering
\begin{tikzpicture}[
    node distance=1.2cm,
    every node/.style={font=\small},
    process/.style={rectangle, rounded corners, draw=navyblue, fill=lightgray, align=center, minimum width=4cm, minimum height=0.9cm},
    decision/.style={diamond, draw=navyblue, fill=lightgray, align=center, aspect=2, inner sep=2pt},
    arrow/.style={-Latex, thick}
]
\node[process] (start) {Start sceny Game};
\node[process, below=of start] (spawn) {BattleSpawner tworzy jednostki};
\node[process, below=of spawn] (queue) {TurnManager buduje kolejkę};
\node[decision, below=of queue] (team) {Czy tura gracza?};
\node[process, below left=1.2cm and 1.0cm of team] (player) {Gracz wybiera ruch/atak};
\node[process, below right=1.2cm and 1.0cm of team] (ai) {AI rozmyte wybiera decyzję};
\node[process, below=3.0cm of team] (combat) {Wykonanie akcji i aktualizacja UI};
\node[decision, below=of combat] (end) {Koniec gry?};
\node[process, below left=1.2cm and 1.0cm of end] (sceneEnd) {Scena Victory/GameOver};
\node[process, below right=1.2cm and 1.0cm of end] (next) {Następna jednostka};

\draw[arrow] (start) -- (spawn);
\draw[arrow] (spawn) -- (queue);
\draw[arrow] (queue) -- (team);
\draw[arrow] (team) -- node[left] {tak} (player);
\draw[arrow] (team) -- node[right] {nie} (ai);
\draw[arrow] (player) -- (combat);
\draw[arrow] (ai) -- (combat);
\draw[arrow] (combat) -- (end);
\draw[arrow] (end) -- node[left] {tak} (sceneEnd);
\draw[arrow] (end) -- node[right] {nie} (next);
\draw[arrow] (next.east) .. controls +(2,0) and +(2,0) .. (team.east);
\end{tikzpicture}
\caption{Ogólny przepływ rozgrywki}
\end{figure}

\newpage

% ============================================================
\section{Sceny gry}
% ============================================================

\subsection{Lista scen}

\begin{longtable}{p{4cm} p{10cm}}
\toprule
\textbf{Scena} & \textbf{Opis} \\
\midrule
\endhead
\texttt{MainMenu} & Ekran startowy gry. Pozwala rozpocząć rozgrywkę albo zamknąć aplikację. \\[4pt]
\texttt{Game} & Główna scena walki. Zawiera planszę, jednostki, tło, kamerę, system tur, UI i log walki. \\[4pt]
\texttt{GameOver} & Scena porażki wyświetlana po pokonaniu wszystkich jednostek gracza. \\[4pt]
\texttt{Victory} & Scena zwycięstwa wyświetlana po pokonaniu wszystkich przeciwników. \\
\bottomrule
\end{longtable}

\subsection{Scena Game}

Najważniejszą sceną projektu jest \texttt{Game}. Znajdują się w niej systemy odpowiedzialne za właściwą rozgrywkę:

\begin{itemize}
    \item obiekty \texttt{Grid} i \texttt{Tilemap} reprezentujące planszę,
    \item \texttt{World\_Background} jako tło dopasowane do widoku kamery,
    \item \texttt{TurnManager} zarządzający kolejką i fazami,
    \item \texttt{BattleSpawner} tworzący jednostki,
    \item \texttt{CombatSystem} obsługujący ataki,
    \item \texttt{FuzzyAiController} sterujący przeciwnikami,
    \item panele UI odpowiadające za przyciski akcji, log, kolejkę i panel celu.
\end{itemize}

\newpage

% ============================================================
\section{Jednostki i klasy postaci}
% ============================================================

\subsection{Model jednostki}

Jednostka jest reprezentowana przez klasę \texttt{Unit}. Obiekt przechowuje najważniejsze informacje potrzebne do walki:

\begin{itemize}
    \item identyfikator jednostki,
    \item nazwę wyświetlaną w interfejsie,
    \item drużynę: gracz albo przeciwnik,
    \item pozycję na siatce,
    \item aktualne i maksymalne punkty życia,
    \item zasięg ruchu,
    \item zasięg i obrażenia ataku wręcz,
    \item zasięg i obrażenia ataku dystansowego,
    \item aktualną postawę bojową.
\end{itemize}

\subsection{Klasy jednostek}

Parametry klas są przechowywane w assetach \texttt{UnitClassData}. Dzięki temu statystyki postaci można zmieniać z poziomu edytora Unity bez edytowania kodu C\#.

\begin{longtable}{p{3.5cm} p{1.8cm} p{1.8cm} p{2cm} p{2cm} p{2cm} p{2cm}}
\toprule
\textbf{Klasa} & \textbf{HP} & \textbf{Ruch} & \textbf{Zasięg wręcz} & \textbf{Zasięg dyst.} & \textbf{Obr. wręcz} & \textbf{Obr. dyst.} \\
\midrule
\endhead
Light Archer & 18 & 5 & 1 & 3 & 5 & 5 \\[4pt]
Light Warrior & 22 & 4 & 1 & 6 & 3 & 4 \\[4pt]
Heavy Archer & 24 & 3 & 1 & 5 & 6 & 7 \\[4pt]
Heavy Warrior & 30 & 3 & 1 & 0 & 9 & 0 \\
\bottomrule
\end{longtable}

\subsection{Drużyny}

W projekcie występują dwie drużyny:

\begin{itemize}
    \item \textbf{Player} --- jednostki kontrolowane przez gracza,
    \item \textbf{Enemy} --- jednostki sterowane przez AI.
\end{itemize}

Obie drużyny korzystają z tych samych klas jednostek, ale posiadają osobne assety konfiguracyjne oraz osobne listy w spawnerze. Dzięki temu można w przyszłości łatwo wprowadzić asymetrię, np. inne statystyki dla przeciwników.

\subsection{Paski życia}

Każda jednostka otrzymuje komponent \texttt{UnitHealthBar}. Komponent tworzy dwa obiekty potomne:

\begin{itemize}
    \item tło paska życia,
    \item wypełnienie paska życia.
\end{itemize}

Wypełnienie jest skalowane proporcjonalnie do aktualnego HP:

\begin{equation}
    \text{procent HP} = \frac{\text{aktualne HP}}{\text{maksymalne HP}}
\end{equation}

Paski gracza są zielone, a paski przeciwników czerwone. Pasek jest umieszczany nad sprite'em postaci i sortowany tak, aby był widoczny nad jednostką.

\newpage

% ============================================================
\section{Plansza i system kafelków}
% ============================================================

\subsection{Reprezentacja planszy}

Plansza jest obsługiwana przez \texttt{GridManager}. Domyślny rozmiar siatki wynosi:

\begin{equation}
    12 \times 12 = 144 \text{ pola}
\end{equation}

Każde pole jest reprezentowane przez \texttt{Vector2Int}, gdzie:

\begin{itemize}
    \item \texttt{x} oznacza kolumnę,
    \item \texttt{y} oznacza wiersz.
\end{itemize}

\subsection{Konwersja współrzędnych}

System \texttt{GridCoordinateSystem} odpowiada za zamianę współrzędnych świata Unity na współrzędne siatki oraz odwrotnie. Jest to potrzebne, ponieważ gracz klika myszą w świecie 2D, a logika gry musi wiedzieć, który kafelek został wskazany.

\begin{longtable}{p{5cm} p{9cm}}
\toprule
\textbf{Operacja} & \textbf{Znaczenie} \\
\midrule
\endhead
\texttt{WorldToGrid} & Zamienia pozycję kliknięcia myszy na współrzędne kafelka. \\[4pt]
\texttt{GridToWorld} & Zamienia współrzędne kafelka na środek pola w świecie Unity. \\
\bottomrule
\end{longtable}

\subsection{Ruch po planszy}

Zasięg ruchu jest obliczany przez przeszukiwanie pól sąsiadujących. Jednostka może wejść tylko na pole:

\begin{itemize}
    \item znajdujące się wewnątrz planszy,
    \item niebędące ścianą ani blokadą,
    \item niezajęte przez inną żywą jednostkę,
    \item osiągalne w ramach jej \texttt{moveRange}.
\end{itemize}

Odległość na planszy jest liczona metryką Manhattan:

\begin{equation}
    d = |x_1 - x_2| + |y_1 - y_2|
\end{equation}

\newpage

% ============================================================
\section{System tur}
% ============================================================

\subsection{Zadania menedżera tur}

Centralnym elementem rozgrywki jest \texttt{TurnManager}. Komponent odpowiada za:

\begin{itemize}
    \item rozpoczęcie bitwy,
    \item utworzenie kolejki żywych jednostek,
    \item wybór aktywnej jednostki,
    \item przełączanie faz tury,
    \item obsługę kliknięć gracza,
    \item uruchamianie tur AI,
    \item czyszczenie podświetleń,
    \item sprawdzanie zwycięstwa i porażki.
\end{itemize}

\subsection{Stany walki}

\begin{longtable}{p{5cm} p{9cm}}
\toprule
\textbf{Stan} & \textbf{Opis} \\
\midrule
\endhead
\texttt{AwaitingTurnStart} & System przygotowuje kolejną turę. \\[4pt]
\texttt{PlayerChooseAction} & Gracz wybiera, czy chce wykonać ruch, atak albo zakończyć turę. \\[4pt]
\texttt{PlayerChooseMove} & Gracz wskazuje pole ruchu. \\[4pt]
\texttt{PlayerChooseAttackTarget} & Gracz wskazuje cel ataku. \\[4pt]
\texttt{ExecutingAiTurn} & Przeciwnik wykonuje decyzję AI. \\
\bottomrule
\end{longtable}

\subsection{Sekwencyjne tury przeciwników}

Przeciwnicy nie wykonują ruchu jednocześnie. Każdy wróg otrzymuje własną turę z kolejki, analizuje sytuację, wykonuje ruch i ewentualny atak, a dopiero potem kolejka przechodzi do następnej jednostki. Dzięki temu rozgrywka jest czytelniejsza, a gracz może obserwować działanie każdej jednostki osobno.

\subsection{Akcje gracza}

W turze gracza możliwe są trzy podstawowe akcje:

\begin{itemize}
    \item \textbf{Move} --- pokazuje pola ruchu i pozwala przemieścić jednostkę,
    \item \textbf{Attack} --- pokazuje możliwe cele i pozwala wykonać atak,
    \item \textbf{End Turn} --- kończy turę aktywnej jednostki.
\end{itemize}

Gracz może także pominąć aktualny krok prawym przyciskiem myszy. Jest to przydatne, gdy jednostka ma wykonać tylko atak bez ruchu albo tylko ruch bez ataku.

\newpage

% ============================================================
\section{System walki}
% ============================================================

\subsection{Walidacja ataku}

Atak jest możliwy tylko wtedy, gdy spełnione są wszystkie warunki:

\begin{itemize}
    \item atakujący i cel istnieją,
    \item obie jednostki są żywe,
    \item jednostki należą do przeciwnych drużyn,
    \item dystans między nimi mieści się w zasięgu ataku wręcz albo dystansowego,
    \item wybrany typ ataku zadaje obrażenia większe od zera.
\end{itemize}

\subsection{Dobór obrażeń}

Jednostka może mieć atak wręcz i atak dystansowy. System wybiera najlepszy dostępny typ ataku dla aktualnego dystansu. Jeśli oba typy są możliwe, używana jest wyższa wartość obrażeń.

\begin{equation}
    \text{obrażenia bazowe} = \max(\text{obrażenia wręcz}, \text{obrażenia dystansowe})
\end{equation}

\subsection{Postawy bojowe}

AI może przypisać jednostce jedną z trzech postaw. Postawa wpływa na obrażenia zadawane oraz otrzymywane.

\begin{longtable}{p{4cm} p{5cm} p{5cm}}
\toprule
\textbf{Postawa} & \textbf{Mnożnik zadawanych obrażeń} & \textbf{Mnożnik otrzymywanych obrażeń} \\
\midrule
\endhead
Defensive & 0{,}85 & 0{,}85 \\[4pt]
Balanced & 1{,}00 & 1{,}00 \\[4pt]
Aggressive & 1{,}15 & 1{,}15 \\
\bottomrule
\end{longtable}

Wartość końcowa obrażeń jest zaokrąglana do liczby całkowitej:

\begin{equation}
    \text{obrażenia końcowe} =
    \operatorname{round}(\text{obrażenia bazowe}
    \times \text{mnożnik atakującego}
    \times \text{mnożnik obrońcy})
\end{equation}

\subsection{Eliminacja jednostki}

Po otrzymaniu obrażeń punkty życia jednostki są zmniejszane, ale nie mogą spaść poniżej zera. Jednostka z \texttt{hp = 0} jest uznawana za pokonaną i zostaje wyłączona w scenie.

\newpage

% ============================================================
\section{Sztuczna inteligencja przeciwników}
% ============================================================

\subsection{Rola AI}

Przeciwnicy są sterowani przez \texttt{FuzzyAiController}. AI nie wykonuje losowego ruchu, lecz analizuje:

\begin{itemize}
    \item własne punkty życia,
    \item dystans do celu,
    \item zagrożenie ze strony jednostek gracza,
    \item możliwe pola ruchu,
    \item możliwość ataku przed ruchem albo po ruchu.
\end{itemize}

Wynikiem pracy AI jest obiekt decyzji zawierający:

\begin{itemize}
    \item wybraną postawę,
    \item wybrane pole docelowe,
    \item identyfikator celu,
    \item informację, czy wykonać atak,
    \item kolejność akcji.
\end{itemize}

\subsection{Zmienne oceniane przez logikę rozmytą}

\begin{longtable}{p{4cm} p{4cm} p{6cm}}
\toprule
\textbf{Zmienna} & \textbf{Zakres interpretacji} & \textbf{Znaczenie} \\
\midrule
\endhead
Poziom HP & niski -- wysoki & Czy jednostka powinna grać ostrożniej. \\[4pt]
Dystans do celu & bliski -- daleki & Czy przeciwnik może szybko zaatakować. \\[4pt]
Zagrożenie & niskie -- wysokie & Ilu graczy może zagrozić jednostce w najbliższym otoczeniu. \\[4pt]
Przetrwanie & słabe -- dobre & Relacja aktualnego HP do potencjalnych obrażeń od graczy. \\
\bottomrule
\end{longtable}

\subsection{Funkcje przynależności}

Do obliczania wartości rozmytych stosowane są funkcje:

\begin{itemize}
    \item trójkątna \texttt{Tri(x, a, b, c)},
    \item trapezoidalna \texttt{Trap(x, a, b, c, d)},
    \item operator \texttt{And} jako minimum,
    \item operator \texttt{Or} jako maksimum.
\end{itemize}

Przykład funkcji trójkątnej:

\begin{equation}
\mu(x)=
\begin{cases}
0, & x \leq a \text{ lub } x \geq c \\
\frac{x-a}{b-a}, & a < x < b \\
1, & x=b \\
\frac{c-x}{c-b}, & b < x < c
\end{cases}
\end{equation}

\subsection{Wybór celu}

AI wybiera cel spośród żywych jednostek gracza. Preferowany jest cel, który:

\begin{enumerate}
    \item może zostać zaatakowany po ruchu,
    \item ma najmniej punktów życia,
    \item znajduje się najbliżej przeciwnika.
\end{enumerate}

Jeśli żaden cel nie może zostać zaatakowany po ruchu, AI wybiera najbliższą jednostkę gracza.

\subsection{Defuzyfikacja postawy}

Postawa przeciwnika jest wyznaczana metodą centroidu dla trzech etykiet:

\begin{itemize}
    \item Defensive,
    \item Balanced,
    \item Aggressive.
\end{itemize}

Ogólna postać wzoru:

\begin{equation}
    y = \frac{\sum(\mu_i \cdot v_i)}{\sum \mu_i}
\end{equation}

gdzie:

\begin{itemize}
    \item $\mu_i$ oznacza stopień aktywacji reguły,
    \item $v_i$ oznacza wartość przypisaną danej postawie.
\end{itemize}

\subsection{Sekwencje akcji AI}

\begin{longtable}{p{4cm} p{10cm}}
\toprule
\textbf{Sekwencja} & \textbf{Opis} \\
\midrule
\endhead
\texttt{MoveOnly} & Przeciwnik tylko przemieszcza się na lepszą pozycję. \\[4pt]
\texttt{AttackOnly} & Przeciwnik atakuje bez zmiany pozycji. \\[4pt]
\texttt{MoveThenAttack} & Przeciwnik najpierw podchodzi, a następnie atakuje. \\[4pt]
\texttt{AttackThenMove} & Przeciwnik najpierw atakuje, a potem zmienia pozycję. \\
\bottomrule
\end{longtable}

\newpage

% ============================================================
\section{Interfejs użytkownika}
% ============================================================

\subsection{Elementy HUD}

\begin{longtable}{p{4cm} p{10cm}}
\toprule
\textbf{Element} & \textbf{Opis} \\
\midrule
\endhead
Panel aktywnej jednostki & Pokazuje, która jednostka aktualnie wykonuje turę. \\[4pt]
Przyciski akcji & Pozwalają wybrać ruch, atak albo zakończenie tury. \\[4pt]
Panel fazy & Informuje, czy gracz wybiera akcję, ruch czy cel ataku. \\[4pt]
Kolejka tur & Pokazuje następne jednostki w kolejności działania. \\[4pt]
Panel celu & Prezentuje informacje o wskazanym przeciwniku. \\[4pt]
Log walki & Wyświetla zdarzenia z walki. \\[4pt]
Paski życia & Pokazują aktualne HP bezpośrednio nad postaciami. \\
\bottomrule
\end{longtable}

\subsection{Tło i widok kamery}

Scena walki korzysta z tła świata \texttt{World\_Background}, które jest renderowane za tilemapą i postaciami. Tło jest dopasowywane do widoku kamery, aby wypełniało cały ekran i nie zostawiało pustych pasków. Tilemapa i jednostki pozostają widoczne nad tłem dzięki ustawieniu warstw sortowania.

\subsection{Czytelność planszy}

Projekt stosuje kilka rozwiązań poprawiających czytelność:

\begin{itemize}
    \item osobne podświetlenie pól ruchu,
    \item osobne podświetlenie pól ataku,
    \item sortowanie sprite'ów według pozycji na osi Y,
    \item paski życia nad każdą postacią,
    \item log komunikatów opisujący wykonywane akcje.
\end{itemize}

\newpage

% ============================================================
\section{Najważniejsze fragmenty implementacji}
% ============================================================

\subsection{Walidacja akcji}

Za sprawdzanie poprawności ruchu i ataku odpowiada \texttt{ActionValidator}. Dzięki temu reguły poprawności akcji są wydzielone z menedżera tur.

\begin{lstlisting}[language={[Sharp]C}, caption={Przykładowa logika walidacji ataku}]
public bool CanAttack(Unit attacker, Unit target)
{
    if (attacker == null || target == null) return false;
    if (!attacker.IsAlive || !target.IsAlive) return false;
    if (attacker.team == target.team) return false;

    int distance = GridManager.Manhattan(attacker.gridPos, target.gridPos);
    return attacker.GetBestDamageAtDistance(distance) > 0;
}
\end{lstlisting}

\subsection{Funkcje przynależności}

Funkcje przynależności są zaimplementowane w osobnej klasie statycznej. Upraszcza to testowanie i ponowne użycie logiki rozmytej.

\begin{lstlisting}[language={[Sharp]C}, caption={Trapezoidalna funkcja przynależności}]
public static float Trap(float x, float a, float b, float c, float d)
{
    if (x <= a || x >= d) return 0f;
    if (x >= b && x <= c) return 1f;
    if (x < b) return (x - a) / (b - a);
    return (d - x) / (d - c);
}
\end{lstlisting}

\subsection{Automatyczne paski życia}

Paski życia są dodawane automatycznie podczas tworzenia jednostek. Dzięki temu nie trzeba ręcznie konfigurować każdego prefabu postaci.

\begin{lstlisting}[language={[Sharp]C}, caption={Dodawanie paska życia do jednostki}]
if (u.GetComponent<UnitHealthBar>() == null)
{
    u.gameObject.AddComponent<UnitHealthBar>();
}
\end{lstlisting}

\newpage

% ============================================================
\section{Zakres projektu}
% ============================================================

\subsection{Co system robi}

\begin{itemize}
    \item Wyświetla menu główne oraz pozwala rozpocząć grę.
    \item Generuje planszę kafelkową 12 na 12 pól.
    \item Tworzy jednostki gracza i przeciwnika na losowych wolnych polach startowych.
    \item Obsługuje klasy postaci z danymi zapisanymi w \texttt{ScriptableObject}.
    \item Prowadzi kolejkę tur dla wszystkich żywych jednostek.
    \item Pozwala graczowi wykonywać ruch i atak.
    \item Steruje przeciwnikami przez AI oparte na logice rozmytej.
    \item Oblicza obrażenia z uwzględnieniem postawy bojowej.
    \item Pokazuje zakres ruchu i ataku na tilemapie.
    \item Wyświetla paski życia nad każdą jednostką.
    \item Pokazuje log walki oraz informacje o aktywnej jednostce.
    \item Wykrywa zwycięstwo i porażkę.
\end{itemize}

\subsection{Czego system nie robi}

\begin{itemize}
    \item Nie posiada trybu wieloosobowego.
    \item Nie zapisuje postępu rozgrywki do pliku.
    \item Nie posiada ekwipunku ani rozwoju postaci.
    \item Nie generuje losowych map z przeszkodami w trakcie gry.
    \item Nie posiada zaawansowanych animacji ataku.
    \item Nie posiada menu ustawień dźwięku, grafiki ani sterowania.
\end{itemize}

\subsection{Możliwe kierunki rozwoju}

\begin{itemize}
    \item Dodanie większej liczby klas postaci.
    \item Dodanie umiejętności specjalnych i efektów statusu.
    \item Dodanie przeszkód terenowych wpływających na ruch.
    \item Dodanie zapisu i odczytu stanu gry.
    \item Dodanie kampanii z kilkoma mapami.
    \item Rozbudowanie AI o pamięć poprzednich tur.
    \item Dodanie animacji ataków i efektów wizualnych.
    \item Dodanie balansu poziomu trudności.
\end{itemize}

\newpage

% ============================================================
\section{Kryteria sukcesu}
% ============================================================

\begin{enumerate}
    \item Gra uruchamia się z menu głównego i przechodzi do sceny walki.
    \item Plansza kafelkowa jest widoczna razem z tłem i postaciami.
    \item Tło pokrywa cały obszar widoczny przez kamerę.
    \item Na planszy pojawiają się jednostki gracza i przeciwników.
    \item Każda jednostka posiada widoczny pasek życia.
    \item Gracz może wybrać ruch aktywnej jednostki.
    \item Gracz może zaatakować przeciwnika znajdującego się w zasięgu.
    \item System blokuje ruch na pola niedostępne i zajęte.
    \item System blokuje atak na jednostki spoza zasięgu.
    \item Przeciwnicy wykonują tury po kolei, a nie wszyscy naraz.
    \item AI wybiera cel oraz wykonuje ruch i atak zgodnie z aktualną sytuacją.
    \item Obrażenia są odejmowane od punktów życia jednostki.
    \item Jednostka z zerowym HP znika z pola walki.
    \item Po pokonaniu wszystkich przeciwników uruchamia się scena zwycięstwa.
    \item Po pokonaniu wszystkich jednostek gracza uruchamia się scena porażki.
\end{enumerate}

\newpage

% ============================================================
\section{Testy funkcjonalne}
% ============================================================

\subsection{Przypadki testowe}

\begin{longtable}{p{1cm} p{4cm} p{5cm} p{4cm}}
\toprule
\textbf{ID} & \textbf{Scenariusz} & \textbf{Kroki} & \textbf{Oczekiwany wynik} \\
\midrule
\endhead
T1 & Uruchomienie gry & Otworzyć scenę \texttt{MainMenu} i kliknąć start. & Gra przechodzi do sceny \texttt{Game}. \\[4pt]
T2 & Widoczność planszy & Uruchomić scenę \texttt{Game}. & Widoczna jest tilemapa, tło i jednostki. \\[4pt]
T3 & Ruch gracza & Wybrać aktywną jednostkę i kliknąć pole w zasięgu. & Jednostka przemieszcza się na wybrane pole. \\[4pt]
T4 & Ruch poza zasięg & Kliknąć pole poza podświetlonym zakresem. & Jednostka nie zmienia pozycji. \\[4pt]
T5 & Atak w zasięgu & Wybrać akcję ataku i kliknąć przeciwnika w zasięgu. & Przeciwnik traci HP, a log pokazuje komunikat. \\[4pt]
T6 & Atak poza zasięgiem & Kliknąć przeciwnika poza zasięgiem. & Atak nie zostaje wykonany. \\[4pt]
T7 & Tura AI & Zakończyć turę gracza, gdy aktywny jest przeciwnik. & Przeciwnik wykonuje decyzję automatycznie. \\[4pt]
T8 & Kolejność AI & Poczekać na kilka tur przeciwników. & Każdy wróg porusza się osobno, kolejno. \\[4pt]
T9 & Paski życia & Zadać obrażenia jednostce. & Pasek życia zmniejsza się proporcjonalnie do HP. \\[4pt]
T10 & Koniec gry & Pokonać wszystkie jednostki jednej drużyny. & Gra przechodzi do sceny zwycięstwa albo porażki. \\
\bottomrule
\end{longtable}

\subsection{Testy techniczne}

\begin{itemize}
    \item Kompilacja projektu C\# powinna kończyć się bez błędów.
    \item W konsoli Unity nie powinny pojawiać się błędy \texttt{NullReferenceException} podczas standardowej rozgrywki.
    \item Wszystkie sceny wskazane w logice gry powinny być dodane do Build Settings.
    \item Assety klas postaci powinny mieć przypisane portrety i wartości statystyk.
    \item Prefab jednostki powinien posiadać komponenty wymagane przez \texttt{BattleSpawner}.
\end{itemize}

\newpage

% ============================================================
\section{Instrukcja uruchomienia}
% ============================================================

\subsection{Wymagania}

\begin{itemize}
    \item Unity 6 lub wersja zgodna z projektem,
    \item zainstalowany moduł obsługi projektów 2D,
    \item środowisko C\# obsługiwane przez Unity,
    \item system Windows, macOS albo Linux z obsługą Unity Editor.
\end{itemize}

\subsection{Uruchomienie w Unity Editor}

\begin{enumerate}
    \item Otworzyć Unity Hub.
    \item Wybrać opcję dodania istniejącego projektu.
    \item Wskazać katalog \texttt{SiamSkirmish}.
    \item Poczekać na import assetów.
    \item Otworzyć scenę \texttt{Assets/Scenes/MainMenu.unity} albo \texttt{Assets/Scenes/Game.unity}.
    \item Kliknąć przycisk Play.
\end{enumerate}

\subsection{Najważniejsze katalogi}

\begin{longtable}{p{5cm} p{9cm}}
\toprule
\textbf{Katalog} & \textbf{Zawartość} \\
\midrule
\endhead
\texttt{Assets/Scripts/AI} & Logika rozmyta i decyzje przeciwników. \\[4pt]
\texttt{Assets/Scripts/Combat} & System walki i obrażeń. \\[4pt]
\texttt{Assets/Scripts/Grid} & Plansza, tilemapa i współrzędne. \\[4pt]
\texttt{Assets/Scripts/Spawning} & Tworzenie jednostek na mapie. \\[4pt]
\texttt{Assets/Scripts/Turns} & Kolejka tur i walidacja akcji. \\[4pt]
\texttt{Assets/Scripts/Units} & Dane jednostek, widok i paski życia. \\[4pt]
\texttt{Assets/Scripts/UI} & Interfejs, menu, log i ekrany końcowe. \\[4pt]
\texttt{Assets/Scenes} & Sceny gry. \\[4pt]
\texttt{Assets/Data/Classes} & Assety klas postaci. \\[4pt]
\texttt{Assets/Sprites} & Grafiki postaci, tła i kafelków. \\
\bottomrule
\end{longtable}

\newpage

% ============================================================
\section{Podsumowanie}
% ============================================================

\textit{Siam Skirmish} realizuje kompletny fundament taktycznej gry turowej RPG. Projekt obejmuje planszę kafelkową, system jednostek, klasy postaci, kolejkę tur, walkę, interfejs użytkownika, warunki zwycięstwa i porażki oraz przeciwników sterowanych przez logikę rozmytą.

Najważniejszym elementem technicznym projektu jest połączenie klasycznego systemu tur z AI, które analizuje sytuację na planszy i wybiera zachowanie na podstawie stopniowych ocen, a nie wyłącznie prostych instrukcji warunkowych. Dzięki temu przeciwnicy mogą reagować na poziom zagrożenia, dystans do celu oraz własne punkty życia.

Projekt ma także przejrzystą strukturę modułową. Osobne skrypty odpowiadają za planszę, walkę, jednostki, AI, UI i tury. Taki podział ułatwia dalszy rozwój gry, testowanie poszczególnych mechanik oraz dodawanie nowych funkcji.

\end{document}
